\section{结果}
\label{sec:results}

%% 结果指南：
%% - 客观呈现发现
%% - 有效使用图表
%% - 适当的统计分析
%% - 与基线或现有方法比较
%% - 突出关键发现

%% 实验设置
\subsection{实验设置}

\subsubsection{数据集描述}
我们使用...评估了我们的方法。

\subsubsection{评估指标}
使用以下指标评估性能：
\begin{itemize}
    \item 准确率: $\text{准确率} = \frac{TP + TN}{TP + TN + FP + FN}$
    \item 精确率: $\text{精确率} = \frac{TP}{TP + FP}$
    \item 召回率: $\text{召回率} = \frac{TP}{TP + FN}$
\end{itemize}

\subsubsection{基线方法}
我们将方法与以下基线进行比较：
\begin{itemize}
    \item 基线1：传统方法 \cite{ref1}
    \item 基线2：最先进方法 \cite{ref2}
    \item 基线3：近期改进方法 \cite{ref3}
\end{itemize}

%% 主要结果
\subsection{性能比较}

表~\ref{tab:results}显示了我们的方法与基线方法的性能比较。

\begin{table}[ht]
\centering
\caption{测试数据集上的性能比较}
\label{tab:results}
\begin{tabular}{lccc}
\hline
方法 & 准确率 (\%) & 精确率 (\%) & 召回率 (\%) \\
\hline
基线1 & 85.2 & 83.1 & 87.3 \\
基线2 & 87.9 & 86.2 & 89.1 \\
基线3 & 89.1 & 88.5 & 90.2 \\
\textbf{我们的方法} & \textbf{92.3} & \textbf{91.7} & \textbf{93.1} \\
\hline
\end{tabular}
\end{table}

%% 详细分析
\subsection{详细分析}

图~\ref{fig:performance}展示了不同参数设置下的性能趋势。

\begin{figure}[ht]
\centering
% \includegraphics[width=0.8\textwidth]{figures/performance_plot.pdf}
\caption{不同参数下的性能分析}
\label{fig:performance}
\end{figure}

%% 统计分析
\subsection{统计分析}

为确保结果的可靠性，我们进行了统计显著性检验...

%% 消融研究
\subsection{消融研究}

我们进行了消融研究以理解不同组件的贡献...